\section{Proofs: automorphism and unitary groups of Cuntz algebras}\label{sec:proofs-cuntz}

This section reproduces, from the preprint of the editor \cite{LupiniCuntz2026} in its version of 12 September 2026, the
proofs of the results stated in Section~\ref{sec:operator-algebras}. The text is that of the
preprint, with British spelling: its sections appear as subsections, its macros
are expanded and its labels renamed, and of its introduction only the
statements of the main results are kept. Citations of the companion preprints
point to their own proof pages where the cited statement is reproduced there.
The preprint has not been refereed. The last subsection derives the statements
of Section~\ref{sec:operator-algebras} from those of the preprint.

\subsection{Main results of the preprint}\label{sec:cu-main-results}

\begin{theorem}\label{thm:cu-main}
\uses{def:ea, thm:ellis}
 Put $D=\mathcal{O}_2\otimes\mathcal{K}$.

\begin{enumerate}
\item The norm group $U_1(D)$ is extremely amenable.

\item For every pure state $\varphi$ of $\mathcal{O}_2$, the point-norm
group 
\begin{equation*}
H_\varphi=\{\alpha\in\mathrm{Aut}(\mathcal{O}_2):\varphi\circ\alpha=\varphi\}
\end{equation*}
is extremely amenable.

\item With point-norm topology on the automorphism group, 
\begin{equation*}
M(\mathrm{Aut}(\mathcal{O}_2))\cong S(\mathcal{O}_2).
\end{equation*}
The state space has its weak-* topology and the action $\alpha\cdot\rho=\rho \circ\alpha^{-1}$.

\item With norm topology on the unitary group, 
\begin{equation*}
M(U(D^+))\cong\mathbb{T}.
\end{equation*}
The action is multiplication through the scalar character of $D^+$.
\end{enumerate}
\end{theorem}

\begin{theorem}\label{thm:cu-finite-main}
 For every finite tuple $u_1,\ldots,u_m\in U_1(D)$
and every $\varepsilon>0$, there are $w_1,\ldots,w_l\in U_1(D)$ and
permutations $\sigma_j$ of $\{1,\ldots,l\}$ such that 
\begin{equation*}
\frac1l\sum_{a=1}^l \left\|w_{\sigma_j(a)}u_j-w_a\right\|<\varepsilon
\qquad(1\leq j\leq m).
\end{equation*}
The witnessing unitaries $w_a$ can be supported on one finite matrix corner
of $D$, with the identity on its complement.
\end{theorem}

\subsection{Dynamical and operator-algebraic preliminaries}

\label{sec:cu-prelim}


\subsubsection{Topologies and compact actions}

All compact spaces are nonempty Hausdorff spaces and all group actions are
jointly continuous. The group $\mathrm{Aut}(B)$ of a separable C*-algebra
is Polish in the point-norm topology. Norm unitary groups and their compact
central quotients are Polish SIN groups. Multiplier unitary groups, when
used, carry the strict topology. These conventions are not interchangeable.

A group is amenable if it has a left-invariant mean on $\mathrm{RUCB}(G)$,
the bounded right-uniformly continuous functions. Equivalently, every
compact action admits an invariant regular probability. A group is extremely
amenable if every compact action has a fixed point.

\begin{definition}\label{def:cu-definition-1}
We call a group \emph{whirly} if every invariant regular probability in
every compact action has support contained in the fixed-point set.
\end{definition}

We use this fixed-support formulation throughout. Whirliness does not assert
existence of invariant probabilities. A whirly amenable group is extremely
amenable. The term is used here for a property of a topological group and
its compact actions; it is an analogue, not a special case, of the
measure-theoretic whirliness of Glasner--Tsirelson--Weiss \cite{GlasnerTsirelsonWeiss2005} for
measure-preserving actions.

\begin{lemma}\label{lem:cu-permanence}
\uses{def:cu-definition-1, def:ea}
 Extreme amenability passes through continuous
homomorphisms with dense image, through finite products, and through groups
with a dense directed union of extremely amenable subgroups. Whirliness
passes through continuous dense homomorphisms and through groups
topologically generated by continuous images of whirly groups.
\end{lemma}

\begin{proof}
Pull back an action along the homomorphism. A point fixed by its image is
fixed by its closure. For products, first take the fixed points of one
factor and then of the other. For a directed union, the closed fixed-point
sets have the finite intersection property. For the whirly statements, apply
the same argument to the closed support of an invariant probability: this
one set is fixed by every generating image and hence by its generated
closure.
\end{proof}

\begin{lemma}\label{lem:cu-kernel}
\uses{def:ea, thm:ellis}
 Let $q:G\to Q$ be a continuous open surjective
homomorphism with extremely amenable kernel. Then $M(G)\cong q^*M(Q)$. For a
compact group $C$ and any topological group $L$, 
\begin{equation*}
M(C\times L)\cong C\times M(L).
\end{equation*}
\end{lemma}

\begin{proof}
In a minimal $G$-flow, the kernel has a nonempty closed fixed-point set.
Normality makes it $G$-invariant, so it is the entire flow. The action
descends continuously through the open quotient. This identifies minimal
flows and their factor maps. For the second assertion, in a minimal $C\times
L$-flow $X$ choose an $L$-minimal subflow $Y$. The set $CY$ is compact and
invariant, hence equals $X$. An $L$-factor map $f:M(L)\to Y$ gives the
factor map $(c,x)\mapsto cf(x)$. The displayed product is minimal under its
independent actions.
\end{proof}

\begin{lemma}[Homogeneous completion]\label{lem:cu-completion-criterion}
\uses{def:coprecompact, def:ea, thm:ellis}
 Let $G$ be a topological group and $H\leq G$
a closed, extremely amenable, co-precompact subgroup. Let $\widehat{G/H}$ be
the completion of $G/H$ in the quotient right uniformity. If $\widehat{G/H}$
is a minimal $G$-flow, then $M(G)\cong\widehat{G/H}$.
\end{lemma}

\begin{proof}
Let $X$ be a compact minimal $G$-flow and choose an $H$-fixed point $x\in X$. Joint continuity of the action, uniformly over the compact space $X$,
makes $g\mapsto gx$ right-uniformly continuous; it is constant on left
cosets of $H$, so it descends to a uniformly continuous map $G/H\to X$ and
extends to $\widehat{G/H}$. The extension is equivariant, and its image is a
nonempty closed invariant subset of $X$, hence all of $X$. Thus $\widehat{G/H}$ factors onto every minimal flow; being minimal itself, it is the
universal one.
\end{proof}

The converse direction is the substantial one and is due to Melleray--Nguyen
Van Th\'e--Tsankov: a Polish group whose universal minimal flow is
metrizable with a $G_\delta$ orbit admits such a subgroup \cite[Theorem 1.2]{MNT2016}.


\subsubsection{Norm transport and finite means}

For a bounded metric $d$, write 
\begin{equation*}
\mathsf{W}_1^d(\mu,\nu)= \inf_{\gamma\in\mathrm{Cpl}(\mu,\nu)} \int
d(x,y)\,d\gamma(x,y).
\end{equation*}
We omit $d$ when it is the operator-norm metric.

\begin{lemma}\label{lem:cu-transport}
 Let $G$ have a bounded compatible bi-invariant metric $d$. Suppose that for every finite $E\subset G$ and $\varepsilon>0$ there is
a finitely supported probability $\mu$ with 
\begin{equation*}
\max_{u\in E}\mathsf{W}_1^d(\mu,(R_u)_*\mu)<\varepsilon, \qquad R_u(x)=xu.
\end{equation*}
Then $G$ is amenable. A compactly supported law with this property may be
replaced by an equal-weight finite law, at arbitrarily small additional
error. For equal-weight laws, permutation couplings attain the minimum
transport cost.
\end{lemma}

\begin{proof}
For a bounded $L$-Lipschitz function $f$, a coupling gives 
\begin{equation*}
\left|\int f(xu)\,d\mu(x)-\int f(x)\,d\mu(x)\right| \leq L\mathsf{W}_1^d(\mu,(R_u)_*\mu).
\end{equation*}
Take a weak-* cluster point of the averaging functionals as $E$ increases
and the error decreases. Bounded Lipschitz functions are uniformly dense in
the bounded uniformly continuous functions on a metric space (apply infimal
convolution to real and imaginary parts). Thus the cluster point is a
right-invariant mean on that algebra. Inversion converts it to a
left-invariant mean; the two uniformities agree by bi-invariance.

For compact support, push the law to a finite $\delta$-net and then
approximate its weights by rationals. The transport error of this
replacement tends to zero. Right translations are isometries, so replacing a
law within transport distance $\tau$ increases each invariance error by at
most $2\tau$. Rational weights give equal weights after repetition. The
coupling polytope for two equal-weight labelled laws is the polytope of
doubly stochastic matrices. Its extreme points are permutation matrices,
which proves the last assertion.
\end{proof}

The full equivalence with finite matching criteria is due to Schneider--Thom 
\cite{SchneiderThom2018}. The elementary implication above is the one needed in the
construction.


\subsubsection{Two CP tools}

We use the following consequences of \cite{KirchbergPhillips2000}.

\begin{lemma}[Finite-target CP repair]\label{lem:cu-repair}
 Let $\mathcal{S}$ be an operator system in a unital C*-algebra $B$, let $F$ be a finite-dimensional C*-algebra, and let $L:\mathcal{S}\to F$ be unital, selfadjoint, and completely bounded. There is a
UCP map $T:B\to F$ such that 
\begin{equation*}
\left\|T|_{\mathcal{S}}-L\right\|_{\mathrm{cb}}\leq\left\|L\right\|_{\mathrm{cb}}-1.
\end{equation*}
\end{lemma}

\begin{proof}
Represent $F$ faithfully and unitally as a block-diagonal subalgebra of $M_d$. Apply \cite[Lemma 1.6]{KirchbergPhillips2000} with target $B(\mathbb{C}^d)=M_d$, and compose
with the UCP block-diagonal conditional expectation onto $F$.
\end{proof}

\begin{lemma}[Two-way CP comparison]\label{lem:cu-kp}
 Let $B_0,B_1$ be separable exact unital C*-algebras, with
unitaries $x_{ij}\in B_i$, $1\leq j\leq m$, and faithful unital embeddings $\iota_i:B_i\to\mathcal{O}_2$. If there are UCP maps $T_i:B_i\to\mathcal{O}_2$
satisfying 
\begin{equation*}
\left\|T_0(x_{0j})-\iota_1(x_{1j})\right\|\leq\eta,\qquad
\left\|T_1(x_{1j})-\iota_0(x_{0j})\right\|\leq\eta,
\end{equation*}
then for every $\tau>0$ there is $z\in U(\mathcal{O}_2)$ such that 
\begin{equation*}
\max_j\left\|z\iota_0(x_{0j})z^*-\iota_1(x_{1j})\right\| <11\sqrt\eta+\tau.
\end{equation*}
The listed unitaries need not generate $B_0$ or $B_1$.
\end{lemma}

\begin{proof}
Apply \cite[Proposition 1.13]{KirchbergPhillipsFields2000} to the discrete two-point continuous
field with fibers $B_0,B_1$, section algebra $B_0\oplus B_1$, and unitary
sections $(x_{0j},x_{1j})$. A two-point discrete space is compact metric,
and exactness of the section algebra is immediate from that of its fibers,
so the hypotheses of \cite[Lemma 1.10]{KirchbergPhillipsFields2000} hold. The directed
comparison quantity $\rho_0$ of \cite[Proposition 1.11]{KirchbergPhillipsFields2000} is an
infimum over unital completely positive maps and, by part~(1) of that
proposition, does not depend on the choice of the injective unital
homomorphisms; the maps $T_0,T_1$ therefore bound both of its values by $\eta $. The conclusion of \cite[Proposition 1.13]{KirchbergPhillipsFields2000} bounds an
infimum over $z\in U(\mathcal{O}_2)$ by $11\sqrt\eta$, whence the
arbitrarily small $\tau$. The proposition is stated for arbitrary continuous
unitary sections and imposes no generating assumption on them. The exponent $1/2$ cannot be improved \cite[Remark 3.11]{KirchbergPhillipsFields2000}.
\end{proof}

We also use the embedding and uniqueness theorems of \cite{KirchbergPhillips2000}: every
separable exact unital algebra embeds unitally in $\mathcal{O}_2$, and any
two such embeddings are approximately unitarily equivalent. All the
assertions with ambient target $\mathcal{O}_2$ also hold in a matrix algebra
over $\mathcal{O}_2$, since $M_k(\mathcal{O}_2)\cong\mathcal{O}_2$. Lemma~\ref{lem:cu-repair}, however, is used with finite-dimensional target; no
CP-extension property of $\mathcal{O}_2$ is being assumed.


\subsection{Boundary factorizations and a complete inverse estimate}

\label{sec:cu-boundary}

Let $\Gamma=\mathbb{F}_2=\langle g_1,g_2\rangle$, let $X=\partial\Gamma$ be
its boundary of infinite reduced words, and put 
\begin{equation*}
\mathcal{E}=C(X)\rtimes_r\Gamma.
\end{equation*}
Write $\lambda_g$ for the canonical unitary and $\alpha_g(f)(\xi)=f(g^{-1}\xi)$. The boundary crossed-product description is classical; see, for
example, Robertson's preprint \emph{Boundary crossed products from the free group} (1990), as cited in~\cite{LupiniCuntz2026}. We give the approximation maps explicitly,
including a proof of the nuclearity needed below.


\subsubsection{Prefix averages}

For $n\geq1$, set $S_n=\{x\in\Gamma:|x|<n\}$ and $r=|S_n|$. Define $m_x\in
C(X)$ by setting $m_x(\xi)=1/n$ if $x$ is one of the first $n$ vertices of
the ray from $1$ to $\xi$, and zero otherwise. The identity vertex is
included. These are locally constant functions and 
\begin{equation}  \label{eq:cu-prefix}
\sum_{x\in S_n}m_x=1,\qquad 0\leq m_x\leq 1/n.
\end{equation}
Define a column over $\mathcal{E}$ by 
\begin{equation*}
a=(a_x)_{x\in S_n},\qquad a_x=\lambda_x^*\sqrt{m_x}.
\end{equation*}
Thus $a^*a=1$. Compression of the regular representation on the Hilbert $C(X) $-module $\ell^2(\Gamma,C(X))$ gives a UCP map 
\begin{equation*}
\theta_n:\mathcal{E}\longrightarrow M_r(C(X)),\qquad
\theta_n(f\lambda_h)_{x,y} =\mathbf{1}_{\{x=hy\}}\alpha_{x^{-1}}(f).
\end{equation*}
The return map $Y\mapsto a^*Ya$ is also UCP.

\begin{lemma}\label{lem:cu-boundary-recovery}
 For $f\in C(X)$ and $h\in\Gamma$, 
\begin{equation*}
a^*\theta_n(f\lambda_h)a=fc_{n,h}\lambda_h,\qquad c_{n,h}=\sum_x\sqrt{m_x\alpha_h(m_{h^{-1}x})},
\end{equation*}
where missing indices have coefficient zero. Moreover, 
\begin{equation*}
0\leq c_{n,h}\leq1,\qquad \left\|1-c_{n,h}\right\|\leq\min\{|h|/n,1\}.
\end{equation*}
Consequently $a^*\theta_n(y)a\to y$ in norm for every $y\in\mathcal{E}$. The
algebra $\mathcal{E}$ is separable, unital and nuclear.
\end{lemma}

\begin{proof}
Matrix multiplication and the covariance relation give the displayed
formula. At $\xi$, the function $c_{n,h}$ is $1/n$ times the size of the
intersection of the first $n$ vertices of the rays from $1$ and from $h$
toward $\xi$. In a tree their common terminal ray starts after at most $|h|$
vertices have been lost from either initial segment. The intersection
therefore has at least $\max\{n-|h|,0\}$ vertices. The upper bound also
follows from Cauchy--Schwarz and \eqref{eq:cu-prefix}. Norm convergence follows
on the algebraic crossed product, and then on $\mathcal{E}$ by contractivity.

For a finite set of inputs and a fixed $n$, a sufficiently fine finite
clopen partition of $X$ gives a finite-dimensional subalgebra $B_n\subset
C(X)$ and a UCP map $J_n:M_r(C(X))\to M_r(B_n)$ approximating their images
under $\theta_n$. Namely, evaluate at one point of each partition atom and
extend constantly on that atom. The maps 
\begin{equation*}
\mathcal{E}\xrightarrow{J_n\theta_n}M_r(B_n) \xrightarrow{Y\mapsto a^*Ya}\mathcal{E }
\end{equation*}
give finite-dimensional completely positive approximations to the identity.
This proves nuclearity. Separability and unitality follow from the
construction.
\end{proof}

We will use embeddings in both directions between $\mathcal{E}$ and $\mathcal{O}_2$. The embedding $\mathcal{E}\hookrightarrow\mathcal{O}_2$
follows from Kirchberg--Phillips \cite{KirchbergPhillips2000}. The other one has a simple
formula. Let $P_x$, for $x=g_1,g_1^{-1},g_2,g_2^{-1}$, be the first-letter
projections, and set $S_x=\lambda_x(1-P_{x^{-1}})$. Then 
\begin{equation*}
S_x^*S_x=1-P_{x^{-1}},\qquad S_xS_x^*=P_x.
\end{equation*}
The elements 
\begin{equation}  \label{eq:cu-cuntz-boundary}
t_1=S_{g_1}+P_{g_1^{-1}},\qquad t_2=S_{g_2}+P_{g_2^{-1}}
\end{equation}
are isometries. Their final projections are $P_{g_1}+P_{g_1^{-1}}$ and $P_{g_2}+P_{g_2^{-1}}$, respectively, and add to $1$. The cross terms vanish
because the initial and final pieces are orthogonal. Thus \eqref{eq:cu-cuntz-boundary} determines a unital embedding $j:\mathcal{O}_2\hookrightarrow\mathcal{E}$, injective by simplicity.


\subsubsection{Dilations and complementary unitaries}

Fix arbitrary $u_1,\ldots,u_m\in U(\mathcal{E})$. Choose the preceding
clopen approximations so that 
\begin{equation*}
t_{n,j}=J_n\theta_n(u_j^*),\qquad
\epsilon_n=\max_j\left\|a^*t_{n,j}a-u_j^*\right\|\longrightarrow0.
\end{equation*}
Each $t_{n,j}$ is a contraction. Its Halmos dilation is 
\begin{equation*}
Q_{n,j}^0= 
\begin{pmatrix}
t_{n,j} & (1-t_{n,j}t_{n,j}^*)^{1/2} \\ 
(1-t_{n,j}^*t_{n,j})^{1/2} & -t_{n,j}^*\end{pmatrix}\in U(M_{2r}(B_n)).
\end{equation*}
Put $b=(a,0)^t$. Then $b^*b=1$, and expansion of the square gives 
\begin{equation}  \label{eq:cu-column-error}
\left\|Q_{n,j}^0bu_j-b\right\|^2\leq2\epsilon_n.
\end{equation}
For $N=2r+1$, the matrix 
\begin{equation*}
V=V_n=\begin{pmatrix}
b & 1-bb^* \\ 
0 & b^*\end{pmatrix}
\in U(M_N(\mathcal{E}))
\end{equation*}
is unitary. Here its row blocks have sizes $2r,1$ and its column blocks have
sizes $1,2r$. In particular its first column is $\widetilde a=(b,0)^t$. Set 
\begin{equation*}
Q_j=Q_{n,j}^0\oplus1,\quad U_j=\mathrm{diag}(u_j,1_{N-1}),\quad
R_j=V^*Q_jVU_j,\quad p=e_{11},\quad q=1-p.
\end{equation*}
With $\delta_n=\sqrt{2\epsilon_n}$, equation \eqref{eq:cu-column-error} says $\left\|(R_j-1)p\right\|\leq\delta_n$.

\begin{lemma}\label{lem:cu-polar}
 If $\delta_n<1$, then $d_j=qR_jq$ is invertible in the $q$-corner. Its polar unitary $c_j$ satisfies 
\begin{equation*}
\left\|c_j-d_j\right\|\leq s_n:=1-\sqrt{1-\delta_n^2},\qquad
\left\|R_j-(p+c_j)\right\|\leq\beta_n:=\sqrt2\delta_n+s_n.
\end{equation*}
\end{lemma}

\begin{proof}
For any unitary $R$, the two operators $(R-1)p$ and $p(R-1)$ have the same
norm: the squares of these norms both equal $\left\|p(2-R-R^*)p\right\|$.
Thus both off-diagonal blocks of $R_j$ have norm at most $\delta_n$. It
follows that 
\begin{equation*}
\left\|q-d_j^*d_j\right\|,\ \left\|q-d_jd_j^*\right\|\leq\delta_n^2.
\end{equation*}
Functional calculus proves invertibility and the asserted polar bound.
Finally, 
\begin{equation*}
R_j-(p+d_j)=(R_j-1)p+pR_jq.
\end{equation*}
The two summands have orthogonal initial supports, and the norm of their sum
is at most the square root of the sum of their squared norms. Adding the
polar error proves the result.
\end{proof}


\subsubsection{The complete inverse}

Consider the finite-dimensional unital algebra and its UCP compression 
\begin{equation}  \label{eq:cu-compression}
F_n=M_{2r}(B_n)\oplus\mathbb{C}\subset M_N(C(X)),\qquad \Phi_n(Y)=qV^*YVq.
\end{equation}
Its range is in $qM_N(\mathcal{E})q$, whose unit is $q$. Since $U_jq=q$, 
\begin{equation}  \label{eq:cu-compression-tuple}
\Phi_n(Q_j)=d_j.
\end{equation}

\begin{theorem}[Complete inverse estimate]\label{thm:cu-inverse}
 For $n>4$, the map $\Phi_n$ is injective and 
\begin{equation*}
\left\|\Phi_n^{-1}:\Phi_n(F_n)\to F_n\right\|_{\mathrm{cb}} \leq
\gamma_n:=(1-2/\sqrt n)^{-1}.
\end{equation*}
The estimate is independent of the clopen partition, $r$, and the matrix
amplification level.
\end{theorem}

\begin{proof}
For $\xi\in X$, the regular covariant representation on $\ell^2(\Gamma)$ is 
\begin{equation*}
\pi_\xi(f)\delta_h=f(h\xi)\delta_h,\qquad
\pi_\xi(\lambda_x)\delta_h=\delta_{xh}.
\end{equation*}
Write $\Pi=\widetilde a\widetilde a^*$. Let $v\in\mathbb{C}^k\otimes\mathbb{C}^N$. The terms of $(1_k\otimes\widetilde a^*)(v\otimes\delta_h)$
corresponding to different $x\in S_n$ have different group coordinates $xh$.
Indeed $a_x^*=\sqrt{m_x}\lambda_x$. Each coefficient has modulus at most $n^{-1/2}$. Orthogonality and the fact that $\widetilde a$ is an isometry
imply 
\begin{equation}  \label{eq:cu-small-projection}
\left\|(1_k\otimes\Pi)(v\otimes\delta_h)\right\| \leq
n^{-1/2}\left\|v\right\|.
\end{equation}

Take $Y\in M_k(M_N(C(X)))$. Choose $\xi\in X$ and a unit vector $v$ with $\left\|Y(\xi)v\right\|=\left\|Y\right\|$; compactness of $X$ and finite
dimensionality permit this. Elements of the coefficient algebra preserve the
group coordinate, so 
\begin{equation*}
\pi_\xi^{(kN)}(Y)(v\otimes\delta_1)=Y(\xi)v\otimes\delta_1.
\end{equation*}
Suppressing matrix amplifications, put $e=1-\Pi$. The identity 
\begin{equation*}
eYe-Y=-\Pi Y-eY\Pi
\end{equation*}
and \eqref{eq:cu-small-projection}, applied to both $v$ and $Y(\xi)v$, give 
\begin{equation*}
\left\|(eYe-Y)(v\otimes\delta_1)\right\| \leq 2n^{-1/2}\left\|Y\right\|.
\end{equation*}
The norm in the crossed product dominates this represented norm. Therefore 
\begin{equation*}
\left\|eYe\right\|\geq(1-2/\sqrt n)\left\|Y\right\|.
\end{equation*}
Since $VqV^*=1-\Pi$, unitary conjugation identifies this compression with
the amplification of $\Phi_n$. Restriction to $M_k(F_n)$, for every $k$,
proves the theorem.
\end{proof}

The only input from the group is the pair of facts that the prefix weights
satisfy $m_x\leq1/n$ and that the translates $xh$, $x\in S_n$, are pairwise
distinct. Freeness enters exactly here: in the tree, the first $n$ vertices
of a ray are distinct group elements, so the column $\widetilde a$ spreads a
unit vector over $\left|S_n\right|$ orthogonal group coordinates with
weights $n^{-1/2}$. No other property of $\mathbb{F}_2$ is used.

The operator system $\Phi_n(F_n)$ has unit $q$, and its inverse map is
unital and selfadjoint. Lemma~\ref{lem:cu-repair} therefore gives a UCP map 
\begin{equation*}
\Psi_n:qM_N(\mathcal{E})q\longrightarrow F_n,\qquad \left\|\Psi_n\Phi_n-\mathrm{id}_{F_n}\right\|_{\mathrm{cb}}\leq\gamma_n-1.
\end{equation*}
Combining this with Lemma~\ref{lem:cu-polar} gives 
\begin{equation}  \label{eq:cu-reverse-comparison}
\left\|\Psi_n(c_j)-Q_j\right\|\leq s_n+\gamma_n-1.
\end{equation}
Thus the inverse estimate produces a reverse CP comparison for the actual
complementary tuple. It is stronger than CP approximation of the original
inputs alone.


\subsection{Simultaneous corrections and norm amenability}

\label{sec:cu-norm}

Let $D=\mathcal{O}_2\otimes\mathcal{K}$, and first suppose the inputs belong
to a copy of $\mathcal{E}$ in a corner $fDf\cong\mathcal{O}_2$. All matrices
from Section~\ref{sec:cu-boundary} are now taken over this corner. The
complementary corner 
\begin{equation*}
C=qM_N(fDf)q\cong\mathcal{O}_2
\end{equation*}
contains $c_1,\ldots,c_m$. Choose a unital embedding $\iota:F_n\to C$. To
apply Lemma~\ref{lem:cu-kp}, take 
\begin{equation*}
B_0=F_n,\qquad B_1=C^*(q,c_1,\ldots,c_m),
\end{equation*}
embed them in $C\cong\mathcal{O}_2$ by $\iota$ and the inclusion, and use
the UCP maps $\Phi_n:B_0\to C$ and $\iota\Psi_n|_{B_1}:B_1\to C$. Equations \eqref{eq:cu-compression-tuple} and \eqref{eq:cu-reverse-comparison} compare the
tuples $(Q_j)$ and $(c_j)$ in both directions, with error at most 
\begin{equation*}
\eta_n=s_n+\gamma_n-1\longrightarrow0.
\end{equation*}
Both source algebras are separable and exact. Lemma~\ref{lem:cu-kp} therefore
gives a unitary $z_n\in U(C)$ such that the representation 
\begin{equation*}
\pi_n=\mathrm{Ad}(z_n)\circ\iota:F_n\to C
\end{equation*}
satisfies, for any chosen sequence $\tau_n>0$ tending to zero, 
\begin{equation}  \label{eq:cu-one-representation}
\max_j\left\|c_j-\pi_n(Q_j)\right\|<11\sqrt{\eta_n}+\tau_n.
\end{equation}
In particular a single representation works for all inputs. It is defined on 
$F_n$, irrespective of whether the chosen tuple generates that algebra.

\begin{proposition}\label{prop:cu-haar-law}
 For each finite tuple of inputs in $U(\mathcal{E})$,
embedded in $fDf$, there are compactly supported probability laws on $U_1(D)$
whose norm transport costs under the corresponding normalised right
translations tend to zero.
\end{proposition}

\begin{proof}
\uses{lem:cu-polar, lem:cu-kp, lem:cu-repair, thm:cu-inverse}
Let $K_n=U(F_n)$, with Haar probability. For $P\in K_n$, put 
\begin{equation*}
\Theta_P=p+\pi_n(P)^*,\qquad w_P=PV_n\Theta_P.
\end{equation*}
These are unitaries in $M_N(fDf)$, with the first expression written in the $p,q$ decomposition. Each $\Theta_P$ commutes with each $U_j=\mathrm{diag}(u_j,1_{N-1})$. Direct multiplication yields 
\begin{equation}  \label{eq:cu-relative}
w_P^*w_{PQ_j}U_j=\Theta_P^*R_j\Theta_{PQ_j}.
\end{equation}
If $R_j$ is replaced by $p+\pi_n(Q_j)$, the right side is the identity,
since its complementary block becomes 
\begin{equation*}
\pi_n(P)\pi_n(Q_j)\pi_n(PQ_j)^*=q.
\end{equation*}
Hence Lemma~\ref{lem:cu-polar} and \eqref{eq:cu-one-representation} imply 
\begin{equation}  \label{eq:cu-haar-bound}
\sup_{P\in K_n}\max_j\left\|w_{PQ_j}U_j-w_P\right\| \leq\beta_n+11\sqrt{\eta_n}+\tau_n\longrightarrow0.
\end{equation}

Realise $M_N(fDf)$ as a finite matrix corner in $D$, with first diagonal
corner $f$. Extend its unitaries by the identity on the complement. Let $\mu_n$ be the push-forward of Haar probability under $P\mapsto w_P$. For
every $j$, the map 
\begin{equation*}
P\longmapsto(w_P,w_{PQ_j}U_j)
\end{equation*}
is a coupling of $\mu_n$ with its right translate by $U_j$. This uses right
invariance of Haar probability on the same compact group $K_n$. Equation \eqref{eq:cu-haar-bound} bounds the transport cost. The image is norm compact
because the parametrization is norm continuous.
\end{proof}

\begin{proof}[Proof of Theorem~\protect\ref{thm:cu-finite-main}]
\proves{thm:cu-finite-main}
\uses{lem:cu-transport, prop:cu-haar-law}
We transfer Proposition~\ref{prop:cu-haar-law} twice. First fix a finite
unitary tuple in a corner $fDf\cong\mathcal{O}_2$. Choose a unital embedding 
$\mathcal{E}\hookrightarrow fDf$, and compose it with $j:\mathcal{O}_2\hookrightarrow\mathcal{E}$ from \eqref{eq:cu-cuntz-boundary}. This composite
embedding and an isomorphism $\mathcal{O}_2\to fDf$ are approximately
unitarily equivalent by \cite[Theorem 1.13]{KirchbergPhillips2000}. Consequently the prescribed
tuple, after one common corner-unitary conjugation, is arbitrarily close to
a tuple in the chosen copy of $\mathcal{E}$. Apply Proposition~\ref{prop:cu-haar-law} and conjugate the resulting law back. The conjugating corner
unitary, extended by $1-f$, belongs to $U_1(D)$. Conjugation is an isometry,
and changing a translating unitary by less than $t$ changes its transport
cost by at most $t$.

Second, let $u_1,\ldots,u_m\in U_1(D)$ be arbitrary. Let $f_l=1_{\mathcal{O}_2}\otimes p_l$, where $p_l$ are increasing finite-rank projections forming
an approximate identity of $\mathcal{K}$. Since $u_j-1\in D$, 
\begin{equation*}
(1-f_l)+f_lu_jf_l\longrightarrow u_j
\end{equation*}
in norm. For sufficiently large $l$, every compression $f_lu_jf_l$ is
invertible in $f_lDf_l$; its polar unitary, extended by $1-f_l$, still
converges to $u_j$. The corner $f_lDf_l\cong M_l(\mathcal{O}_2)\cong\mathcal{O}_2$ is of the kind already treated. Disjoint rank-$l$ blocks in $\mathcal{K}$ supply as many equivalent diagonal corners as the construction requires.
The first transfer, followed by this truncation, gives laws with arbitrarily
small costs for all prescribed inputs.

Finally apply Lemma~\ref{lem:cu-transport}: compactly supported laws can be
approximated in norm transport by equal-weight finite laws. Right
translations are isometries, so an approximation within $t$ adds at most $2t$
to every cost. For equal-weight atomic laws, a minimum-cost coupling is a
permutation coupling by the doubly stochastic matrix theorem. All sampled
unitaries lie in the same finite corner used in the construction. This is
exactly the assertion of Theorem~\ref{thm:cu-finite-main}.
\end{proof}

\begin{corollary}\label{cor:cu-strong-diagonal}
 Let $D=\mathcal{O}_2\otimes\mathcal{K}$. For
every finite set $E\subset U(D^+)$ and every $\eta>0$, there are $z_1,\ldots,z_M\in U_1(D)$, supported on one finite matrix corner, such that 
\begin{equation*}
\Delta=\frac1M\sum_{t=1}^M z_t^*\otimes z_t \in D^+\widehat\otimes_\pi D^+
\end{equation*}
lies in $\mathrm{co}\{v^*\otimes v:v\in U_1(D)\}$ and satisfies 
\begin{equation*}
m(\Delta)=1,\qquad \left\|\Delta\right\|_\pi\leq1,\qquad \mathrm{flip}(\Delta)=\Delta,\qquad \max_{g\in E}\left\|g\Delta-\Delta g\right\|_\pi<\eta,
\end{equation*}
where $m(x\otimes y)=xy$ and the tensor product has its usual outer bimodule
actions. Consequently $D^+$ has a strong approximate diagonal with
finite-corner witnesses.
\end{corollary}

\begin{proof}
\uses{thm:cu-finite-main}
Write every $g\in E$ as $g=\zeta_g u_g$, where $\zeta_g\in\mathbb{T}$ and $u_g\in U_1(D)$. Apply Theorem~\ref{thm:cu-finite-main} to the finite tuple $(u_g^*)_{g\in E}$, with tolerance $\eta/2$. Thus there are $w_1,\ldots,w_l$,
supported on one finite corner, and permutations $\sigma_g$ such that 
\begin{equation*}
\frac1l\sum_{b=1}^l \left\|w_{\sigma_g(b)}u_g^*-w_b\right\|<\frac\eta2.
\end{equation*}
Index a new list, counting repetitions, by ordered pairs and put $z_{a,b}=w_a^*w_b$, so that $M=l^2$. It stays in the same corner and is
invariant under taking adjoints, since $z_{a,b}^*=z_{b,a}$. Hence the
element 
\begin{equation*}
\Delta=\frac1{l^2}\sum_{a,b=1}^l z_{a,b}^*\otimes z_{a,b}
\end{equation*}
is flip invariant; also $m(\Delta)=1$ and $\left\|\Delta\right\|_\pi\leq1$.

Fix $g\in E$, and reindex the sum by $\rho_g(a,b)=(a,\sigma_g(b))$. If 
\begin{equation*}
e_{a,b}=\left\|z_{\rho_g(a,b)}u_g^*-z_{a,b}\right\|,
\end{equation*}
then $e_{a,b}=\left\|w_{\sigma_g(b)}u_g^*-w_b\right\|$. The two factors in 
\begin{equation*}
u_g z_{\rho_g(a,b)}^*\otimes z_{\rho_g(a,b)} \quad\hbox{and}\quad
z_{a,b}^*\otimes z_{a,b}u_g
\end{equation*}
differ by at most $e_{a,b}$ each. The projective tensor estimate and the
displayed average therefore give 
\begin{equation*}
\left\|u_g\Delta-\Delta u_g\right\|_\pi \leq\frac2{l^2}\sum_{a,b=1}^l
e_{a,b}<\eta.
\end{equation*}
The scalar $\zeta_g$ does not change this bound. Every element of the unital 
C*-algebra $D^+$ is a linear combination of four unitaries. Directing
these elements by finite subsets and decreasing $\eta$ therefore gives
centrality on all of $D^+$, and hence the asserted strong approximate
diagonal.
\end{proof}

\begin{remark}\label{rem:cu-remark-1}
\uses{thm:cu-finite-main}
The qualitative conclusion also follows from norm amenability through \cite[Theorem 2.4]{AlekseevSchmidtThom2025}. The construction above retains the quantitative
centrality bound and the common finite-corner support furnished by Theorem~\ref{thm:cu-finite-main}.
\end{remark}

\begin{corollary}\label{cor:cu-w-amenable}
 The norm group $U_1(\mathcal{O}_2\otimes\mathcal{K})$
is amenable.
\end{corollary}

\begin{proof}
\uses{lem:cu-transport, thm:cu-finite-main}
Apply the transport criterion in Lemma~\ref{lem:cu-transport} to Theorem~\ref{thm:cu-finite-main}.
\end{proof}

\begin{remark}\label{rem:cu-remark-2}
\uses{prop:cu-haar-law}
The corrections above take place in the full $\mathcal{O}_2$ tail corner.
They are not required to remain in the boundary coefficient algebra. This
freedom is essential: the argument uses Kirchberg--Phillips uniqueness after
the reverse CP map has been constructed. Neither a column averaging
statement alone nor independent corrections for individual generators would
establish the common-law assertion \eqref{eq:cu-haar-bound}.
\end{remark}


\subsection{Extreme amenability and pure-state stabilisers}

\label{sec:cu-stabilizers}


\subsubsection{Whirliness of stable identity components}

We first record the concentration input in a form that also applies to the
general Kirchberg discussion.

\begin{lemma}\label{lem:cu-stable-whirly}
 Let $B$ be a unital C*-algebra such that $B\otimes\mathcal{K}$ has real rank zero. The identity component $U_1(B\otimes \mathcal{K})_0$, with its norm topology, is whirly.
\end{lemma}

\begin{proof}
\uses{lem:cu-permanence}
The Hilbert--Schmidt unitary group 
\begin{equation*}
U_2=\{u\in U(\ell^2):u-1\text{ is Hilbert--Schmidt}\}
\end{equation*}
with its unnormalized Hilbert--Schmidt metric is a L\'evy group \cite[Example 2.15]{GiordanoPestov2007}. Its invariant probabilities on compact flows are
supported on fixed points, by \cite[Theorem 1.1 and Remark 1.4]{GlasnerTsirelsonWeiss2005}. The
inclusion into norm $U_1(\mathcal{K})$ is continuous and has dense image,
since finite-rank perturbations of the identity are dense. Thus norm $U_1(\mathcal{K})$ is whirly by Lemma~\ref{lem:cu-permanence}.

Every nonzero projection $p\in B\otimes\mathcal{K}$ can be the first
diagonal matrix unit of a copy of $\mathcal{K}$ in $B\otimes\mathcal{K}$. To
see this, compress $p$ to a sufficiently large canonical finite corner and
use functional calculus to obtain a projection $p^{\prime }$ with $\left\|p-p^{\prime }\right\|<1$. A unitary in the unitisation conjugates $p^{\prime }$ to $p$. Infinitely many disjoint copies of the finite corner
give equivalent copies of $p^{\prime }$, and partial isometries between them
give matrix units. Conjugate these back; the resulting matrix units remain
in the ideal $B\otimes\mathcal{K}$, since a unitary of its unitisation
normalises that ideal. In particular the projection circle 
\begin{equation*}
\{1+(z-1)p:z\in\mathbb{T}\}
\end{equation*}
is contained in a copy of $U_1(\mathcal{K})$.

By real rank zero, every selfadjoint element of $B\otimes\mathcal{K}$ is a
norm limit of finite-spectrum selfadjoint elements. After zero spectral
projections are omitted, their exponentials are products of
projection-circle elements. Such exponentials generate $U_1(B\otimes\mathcal{K})_0$. If a compact action of this group has an invariant probability, its
support is fixed by each of the copies of $U_1(\mathcal{K})$, then by the
group they generate, and finally by its closure. This proves whirliness.
\end{proof}

\begin{theorem}\label{thm:cu-w-ea}
 The norm group $W=U_1(\mathcal{O}_2\otimes\mathcal{K})$ is
extremely amenable.
\end{theorem}

\begin{proof}
\uses{cor:cu-w-amenable, lem:cu-stable-whirly}
The algebra $\mathcal{O}_2\otimes\mathcal{K}$ has real rank zero, and its $K_1$-group vanishes \cite[Theorem 1.2(ii)]{Zhang1992}, \cite{Cuntz1981}. For this
stabilised algebra, stability identifies $W/W_0$ with $K_1(\mathcal{O}_2\otimes\mathcal{K})$ \cite[8.1.1]{Blackadar1998}, so $W=W_0$. Lemma~\ref{lem:cu-stable-whirly} makes $W$ whirly. Corollary~\ref{cor:cu-w-amenable}
supplies an invariant probability in every compact $W$-flow, and its
nonempty support consists of fixed points. Thus $W$ is extremely amenable.
\end{proof}


\subsubsection{The hereditary algebra of a pure state}

Fix a pure state $\varphi$ on $A=\mathcal{O}_2$, with GNS triple $(\pi_\varphi,\mathcal{H}_\varphi,\xi_\varphi)$, and define 
\begin{equation*}
D_\varphi=\{a\in A: \varphi(a^*a)=\varphi(aa^*)=0\}.
\end{equation*}
This is the hereditary subalgebra with open support $1-p_\varphi$, where $p_\varphi$ is the minimal support projection of $\varphi$ in $A^{**}$.

\begin{lemma}\label{lem:cu-hereditary}
 The algebra $D_\varphi$ is isomorphic to $\mathcal{O}_2\otimes\mathcal{K}$. Under its unitisation embedding in $A$, 
\begin{equation*}
U_1(D_\varphi) =\{u\in U(A):\pi_\varphi(u)\xi_\varphi=\xi_\varphi\}.
\end{equation*}
\end{lemma}

\begin{proof}
The hereditary algebra is nonzero: otherwise its support would vanish and $1=p_\varphi$ would be minimal, forcing $A=\mathbb{C}$. It is full by
simplicity. It is nonunital, since a unit for it would equal $1-p_\varphi$
and would put the nonzero minimal projection $p_\varphi$ in $A$, impossible
in a purely infinite simple algebra. It is separable. Zhang's stability
theorem \cite[Theorem 1.2(i)]{Zhang1992} makes this nonunital hereditary algebra
stable, and Brown's full hereditary subalgebra theorem \cite[Theorem 2.8]{Brown1977} identifies its stabilisation with $A\otimes\mathcal{K}$. Thus $D_\varphi\cong A\otimes\mathcal{K}$.

For a unitary $u$, the condition $\pi_\varphi(u)\xi_\varphi =\xi_\varphi$
implies the same condition for $u^*$. It is therefore equivalent to $u-1\in
D_\varphi$, directly from the definition.
\end{proof}

\begin{proposition}\label{prop:cu-stabilizer-density}
 The group $H_\varphi=\mathrm{Aut}(A,\varphi)$
is a closed Polish subgroup of $\mathrm{Aut}(A)$, and the continuous
homomorphism 
\begin{equation*}
\mathrm{Ad}:U_1(D_\varphi)\longrightarrow H_\varphi
\end{equation*}
has dense image.
\end{proposition}

\begin{proof}
Closedness follows from continuity of $\alpha\mapsto\varphi\circ\alpha$ for
the point-norm and weak-* topologies. Thus $H_\varphi$ is Polish. Let $\alpha\in H_\varphi$, and prescribe finite point-norm tests and an error. By 
$\mathcal{O}_2$-uniqueness \cite[Theorem 1.13]{KirchbergPhillips2000}, approximate $\alpha$ on
these tests and on an additional finite state-correction test by $\gamma=\mathrm{Ad} u$, $u\in U(A)$. The state $\psi=\varphi\gamma$ is unitarily
equivalent to $\varphi$ and sufficiently close to it on the additional test,
since $\varphi\alpha=\varphi$.

Property 1.2 of \cite{KishimotoOzawaSakai2003} gives a continuous path of unitaries from the
identity implementing $\psi=\varphi\mathrm{Ad} v$ with $\mathrm{Ad} v$ as
close to the identity on the original tests as prescribed. Applicability
uses that an irreducible representation of $A$ contains no nonzero compact
operators. Nuclearity of $A$ is not required for Property~1.2: \cite[Theorem~4.3]{KishimotoOzawaSakai2003} proves Property~1.3 for every C*-algebra, and \cite[Theorem~3.1]{KishimotoOzawaSakai2003} upgrades Property~1.3 to Property~1.2. Put $z=uv^*$. Then $\varphi\mathrm{Ad} z=\varphi$, and 
\begin{equation*}
\left\|\mathrm{Ad} z(a)-\alpha(a)\right\| \leq\left\|\mathrm{Ad}
v^*(a)-a\right\|+\left\|\gamma(a)-\alpha(a)\right\|.
\end{equation*}
Both terms are small. Equality of the two pure vector states implies $\pi_\varphi(z)\xi_\varphi=\lambda\xi_\varphi$ for some $\lambda\in\mathbb{T}$. Replacing $z$ by $\lambda^{-1}z$ does not change its adjoint action and
puts it in $U_1(D_\varphi)$. Finally, norm continuity of the adjoint map
follows from $\left\|\mathrm{Ad} u(a)-\mathrm{Ad} v(a)\right\|\leq2\left\|a\right\|\left\|u-v\right\|$.
\end{proof}

\begin{corollary}\label{cor:cu-h-ea}
 Every pure-state stabiliser $H_\varphi$ in $\mathrm{Aut}(\mathcal{O}_2)$ is extremely amenable.
\end{corollary}

\begin{proof}
\uses{lem:cu-hereditary, prop:cu-stabilizer-density, thm:cu-w-ea}
Apply Theorem~\ref{thm:cu-w-ea}, Lemma~\ref{lem:cu-hereditary}, Proposition~\ref{prop:cu-stabilizer-density}, and permanence under continuous homomorphisms
with dense image.
\end{proof}

\begin{corollary}\label{cor:cu-scalar-umf}
 The universal minimal flow of norm $U((\mathcal{O}_2\otimes\mathcal{K})^+)$ is $\mathbb{T}$, acted on through the scalar
quotient.
\end{corollary}

\begin{proof}
\uses{lem:cu-kernel}
The scalar character $\chi:D^+\to\mathbb{C}$ gives the topological direct
product 
\begin{equation*}
U(D^+)\longrightarrow\mathbb{T}\times W,\qquad u\longmapsto\bigl(\chi(u),\overline{\chi(u)}u\bigr).
\end{equation*}
The normal subgroup $W$ is extremely amenable. Apply Lemma~\ref{lem:cu-kernel};
the UMF of the compact quotient $\mathbb{T}$ is its translation action.
\end{proof}


\subsection{Uniform pure-state homogeneity and the state-space flow}

\label{sec:cu-state}

The local argument in this section is stated for a larger class of algebras.
Its purpose is to make the uniformity in the base state explicit. All state
neighbourhoods below are weak-* neighbourhoods.


\subsubsection{Approximately central UCP rows}

\begin{lemma}\label{lem:cu-central-row}
 Let $A$ be unital and nuclear. For every finite $F\subset A$ and $\tau>0$, there are $x_1,\ldots,x_n\in A$ such that 
\begin{equation*}
\sum_i x_ix_i^*=1,\qquad
\sup_{\left\|b\right\|\leq1}\left\|[a,\Phi(b)]\right\|<\tau\quad(a\in F),
\qquad \Phi(b)=\sum_i x_i b x_i^*.
\end{equation*}
The row depends only on $F,\tau$.
\end{lemma}

\begin{proof}
Apply \cite[Lemma 4.2]{KishimotoOzawaSakai2003} to the universal representation. Nuclearity
makes its weak closure injective. The lemma gives finite positive-form
tensors 
\begin{equation*}
T=\sum_i y_i\otimes y_i^*,\qquad \sum_i\left\|y_i\right\|^2\leq1,
\end{equation*}
with $\left\|aT-Ta\right\|_\pi$ arbitrarily small on $F$, and with $c=\sum_i
y_iy_i^*$ converging weakly to $1$. Here $\left\|\cdot\right\|_\pi$ is the
Banach projective tensor norm, with the usual outer bimodule actions.
Convexification of tails makes $c\to1$ in norm, preserving the positive form
and the centrality estimates: concatenate rows multiplied by square roots of
the convex weights. Also $0\leq c\leq1$. Choose $c$ invertible and put $d=c^{-1/2}$, $x_i=dy_i$.

For $\Psi(b)=\sum_i y_i b y_i^*$, the projective estimate gives $\left\|\mathrm{ad}(a)\circ\Psi\right\|\leq\left\|aT-Ta\right\|_\pi$. Moreover 
\begin{equation*}
\left\|\Phi-\Psi\right\|\leq(\left\|d\right\|+1)\left\|d-1\right\|.
\end{equation*}
Consequently 
\begin{equation*}
\left\|\mathrm{ad}(a)\circ\Phi\right\| \leq\left\|aT-Ta\right\|_\pi+
2\left\|a\right\|(\left\|d\right\|+1)\left\|d-1\right\|.
\end{equation*}
Both terms can be made arbitrarily small, and $\sum_i x_ix_i^*=dcd=1$.
\end{proof}

\begin{lemma}[Gram comparison]\label{lem:cu-gram}
 Let $X,Y:\mathbb{C}^n\to\mathcal{H}$ have orthogonal
ranges. If their Gram matrices $G_0=X^*X$, $G_1=Y^*Y$ differ entrywise by at
most $\delta$, there is a positive contraction $f$ on the sum of their
ranges such that, on each standard basis vector $e_i$, 
\begin{equation*}
\left\|fXe_i-(Xe_i-Ye_i)/2\right\|,\ \left\|fYe_i-(Ye_i-Xe_i)/2\right\|
\leq\tfrac12\sqrt{n\delta}.
\end{equation*}
\end{lemma}

\begin{proof}
Write $X=U\sqrt{G_0}$, $Y=V\sqrt{G_1}$ in polar form and put $Z=VU^*$. Since
the ranges are orthogonal, $Z+Z^*$ is a selfadjoint contraction on their
sum. Set $f=(1-Z-Z^*)/2$. The support identities give 
\begin{equation*}
ZX-Y=V(\sqrt{G_0}-\sqrt{G_1}),\qquad Z^*Y-X=U(\sqrt{G_1}-\sqrt{G_0}).
\end{equation*}
Also $Z^*X=ZY=0$. The square-root estimate $\left\|\sqrt{G_0}-\sqrt{G_1}\right\|\leq\left\|G_0-G_1\right\|^{1/2} \leq\sqrt{n\delta}$ proves the
claim, including singular Gram matrices.
\end{proof}

\begin{lemma}[Glimm approximation with avoidance]\label{lem:cu-glimm-avoid}
 Let $A$ be separable and unital, and let $\pi:A\to B(\mathcal{H})$ be a faithful irreducible representation with $\pi(A)\cap \mathcal{K}(\mathcal{H})=0$. Given a state $\theta$, a finite set $E\subset
A $, a finite-dimensional subspace $M\subset\mathcal{H}$, and $r>0$, there
is a unit vector $\zeta\in M^\perp$ such that 
\begin{equation*}
\bigl|\langle\pi(a)\zeta,\zeta\rangle-\theta(a)\bigr|<r \qquad(a\in E).
\end{equation*}
\end{lemma}

\begin{proof}
Glimm's vector-state approximation theorem \cite{Glimm1961} gives unit vectors $\zeta_j$ converging weakly to zero whose vector states converge weak-* to $\theta$. If $P_M$ is the projection onto $M$, then $\left\|P_M\zeta_j\right \|\to0$. Consequently the normalised vectors 
\begin{equation*}
\frac{(1-P_M)\zeta_j}{\left\|(1-P_M)\zeta_j\right\|}
\end{equation*}
belong to $M^\perp$, converge in norm to $\zeta_j$, and retain the required
finite weak-* approximation. Here one may use 
\begin{equation*}
\bigl|\langle\pi(a)\xi,\xi\rangle- \langle\pi(a)\eta,\eta\rangle\bigr| \leq
2\left\|a\right\|\left\|\xi-\eta\right\|
\end{equation*}
for unit vectors $\xi,\eta$.
\end{proof}


\subsubsection{Uniform local correction}

\begin{proposition}\label{prop:cu-approx-alignment}
 Let $A$ be simple, separable, unital, infinite
dimensional and nuclear. Given finite $F\subset A$ and $\varepsilon>0$,
there are finite $E\subset A$ and $\delta>0$ with the following property. If
pure states $\omega,\nu$ satisfy $\left|\omega(a)-\nu(a)\right|<\delta$ for $a\in E$, then for every finite $K\subset A$ and $t>0$ there is a unitary
path $(u_s)_{0\leq s\leq1}$, starting at $1$, such that 
\begin{equation*}
\sup_s\max_{a\in F}\left\|\mathrm{Ad} u_s(a)-a\right\|<\varepsilon,\qquad
\max_{a\in K}\left|\omega(\mathrm{Ad} u_1(a))-\nu(a)\right|<t.
\end{equation*}
If the GNS representations of the two pure states are unitarily equivalent, $u_1$ may be chosen so that $\omega\circ\mathrm{Ad} u_1=\nu$.
\end{proposition}

\begin{proof}
\uses{lem:cu-central-row, lem:cu-glimm-avoid, lem:cu-gram}
Let $c=\max\bigl(\{1\}\cup\{\left\|a\right\|:a\in F\}\bigr)$. It suffices to
prove the assertion for $c^{-1}F$ with tolerance $c^{-1}\varepsilon$;
relabeling, assume that $F$ is in the unit ball. Choose $\tau,\kappa>0$ such
that 
\begin{equation*}
\kappa<1/(16\pi),\qquad 2\pi\tau+40\pi\kappa<\varepsilon.
\end{equation*}
Take the row of Lemma~\ref{lem:cu-central-row} with commutator bound $\tau$,
let its length be $n$, and put 
\begin{equation*}
E=\{x_ix_j^*:1\leq i,j\leq n\},\qquad \delta_0=\kappa^2/(2n^2).
\end{equation*}
These choices precede the choice of either state.

First suppose that $\omega,\nu$ are represented by unit vectors $\xi,\eta$
in one irreducible representation $\pi$, and are $\delta_0$-close on $E$.
The representation is faithful and essential: a nonzero intersection of $\pi(A)$ with the compact operators would, by simplicity, make $A$
elementary. Lemma~\ref{lem:cu-glimm-avoid}, applied to the state of $\xi$,
gives a unit vector $\zeta$ whose state is $\delta_0$-close to that of $\xi$
on $E$, and which is orthogonal to 
\begin{equation*}
\mathrm{span}\{\pi(x_ix_j^*)\xi,\pi(x_ix_j^*)\eta:1\leq i,j\leq n\}.
\end{equation*}
The displayed span contains $\xi,\eta$, because $\sum_i x_ix_i^*=1$. The
row-adjoint spans belonging to $\zeta$ and to either $\xi$ or $\eta$ are
orthogonal. The corresponding Gram matrices differ entrywise by less than $2\delta_0$.

Apply Lemma~\ref{lem:cu-gram} to the tuples $(\pi(x_i^*)\xi)_i$ and $(\pi(x_i^*)\zeta)_i$. Let $M$ be their finite-dimensional summed span.
Extend the positive contraction furnished by Lemma~\ref{lem:cu-gram} by zero on 
$M^\perp$. Positive Kadison transitivity \cite[Theorem II.6.1.13]{Blackadar2006} gives $a\in A_+$ whose restriction to $M$ is that contraction.
The subspace $M$ reduces $\pi(a)$, so replacing $a$ by $h=\min(a,1)$ through
continuous functional calculus preserves its restriction and gives $0\leq
h\leq1$. Put $k=\Phi(h)$. Multiplication by the row $(\pi(x_i))_i$, whose
norm is one, shows 
\begin{equation*}
\left\|\pi(k)\xi-(\xi-\zeta)/2\right\|<\kappa,\qquad
\left\|\pi(k)\zeta-(\zeta-\xi)/2\right\|<\kappa.
\end{equation*}
Indeed the bound before enlarging it to $\kappa$ is $\sqrt n\,\sqrt{2n\delta_0}/2=\kappa/2$. Thus 
\begin{equation*}
\left\|\pi(k)(\xi+\zeta)\right\|<2\kappa,\qquad
\left\|(\pi(k)-1)(\xi-\zeta)\right\|<2\kappa.
\end{equation*}
The integral formula for the exponential gives 
\begin{equation*}
\begin{split}
\left\|(\pi(e^{i\pi k})-1)(\xi+\zeta)\right\|&<2\pi\kappa, \\
\left\|(\pi(e^{i\pi k})+1)(\xi-\zeta)\right\|&<2\pi\kappa.
\end{split}\end{equation*}
Taking their half-sum yields 
\begin{equation*}
\left\|\pi(e^{i\pi k})\xi-\zeta\right\|<2\pi\kappa.
\end{equation*}
The path $e^{is\pi k}$ moves each $a\in F$ by less than $\pi\tau$.

For completeness, two unit vectors at distance $d<1/4$ can be joined by a
unitary on their two-dimensional span whose selfadjoint logarithm has norm
at most $4d$. Extend the two-vector isometry to that span and use the
principal logarithm; the estimate also follows from the elementary plane
rotation formula. Norm-controlled selfadjoint Kadison transitivity \cite[Theorem II.6.1.13]{Blackadar2006} lifts the logarithm to an element of $A$ of
norm less than $5d$. Its exponential path therefore has norm distance from $1 $ less than $5d$, and moves each element of the unit ball by less than $10d<20\pi\kappa$. Apply this to correct $\pi(e^{i\pi k})\xi$ exactly to $\zeta$. The concatenated path ends at $U$ with $\pi(U)\xi=\zeta$ and moves $F $ by less than 
\begin{equation*}
B=\pi\tau+20\pi\kappa.
\end{equation*}
The same construction for $\eta,\zeta$ gives $V$ with $\pi(V)\eta=\zeta$ and
the same bound. Concatenate the reversed adjoint path $U_s^*$ with $U^*V_s$.
Its endpoint $w=U^*V$ satisfies $\pi(w^*)\xi=\eta$, so $\omega\mathrm{Ad}
w=\nu $, and its movement on $F$ is less than $2B<\varepsilon$.

For arbitrary pure states take $\delta=\delta_0/3$. In the essential
irreducible representation of $\omega$, Lemma~\ref{lem:cu-glimm-avoid} (with no
avoidance subspace) approximates $\nu$ on $E\cup K$ by a vector state $\nu^{\prime }$, with error smaller than $\min\{\delta_0/3,t\}$. The vector
state is pure and has GNS representation equivalent to that of $\omega$,
since $\pi$ is irreducible. Moreover $\omega,\nu^{\prime }$ are less than $2\delta_0/3$, hence $\delta_0$, apart on $E$. Apply the exact
equivalent-state result just proved.
\end{proof}

\begin{theorem}[Uniform exact alignment]\label{thm:cu-alignment}
\uses{prop:cu-approx-alignment}
 Under the hypotheses of Proposition~\ref{prop:cu-approx-alignment}, for every finite $F\subset A$ and $\varepsilon>0$
there are finite $E\subset A$ and $\delta>0$ such that any two pure states $\rho_0,\rho_1$ that are $\delta$-close on $E$ satisfy 
\begin{equation*}
\rho_1=\rho_0\circ v^{-1},\qquad \max_{a\in
F}\left\|v(a)-a\right\|<\varepsilon
\end{equation*}
for some $v\in\mathrm{Aut}(A)$.
\end{theorem}

\begin{proof}
\uses{prop:cu-approx-alignment}
We give the alternating construction, including inverse control. We write $\rho\gamma=\rho\circ\gamma$; thus $(\rho\mathrm{Ad} R)\mathrm{Ad} u=\rho \mathrm{Ad}(Ru)$. Let $(a_j)$ be dense in the unit ball and choose positive $(\eta_n)_{n\geq1}$ with $\sum_n\eta_n<\varepsilon/4$. Set $P=Q=1$ initially
and $L_0=F\cup\{a_1\}$. Let $(G_0,\delta_0)$ be the test from Proposition~\ref{prop:cu-approx-alignment} for $(L_0,\eta_1)$. This is the required initial
test, independent of the states. Assume the states pass it.

At stage $n\geq1$, update $Q$ if $n$ is odd and $P$ if $n$ is even. Denote
the product to be updated by $R$ and the other one by $S$. Before choosing
its correction, set 
\begin{equation*}
L_n=F\cup\mathrm{Ad}(S^*)(F)\cup \{a_j,\mathrm{Ad}(S^*)(a_j):j\leq n+1\}.
\end{equation*}
Choose the next test $(G_n,\delta_n)$ for control $(L_n,\eta_{n+1})$,
enlarge $G_n$ by $G_{n-1}$ and $\{a_1,\ldots,a_n\}$, and decrease $\delta_n$
below $2^{-n}$. The current states $\rho_0\mathrm{Ad} Q,\rho_1\mathrm{Ad} P$
are $\delta_{n-1}$-close on $G_{n-1}$. Apply the proposition with control $(L_{n-1},\eta_n)$, ordering the two states so that the active one is
changed, and request final closeness on $G_n$ within $\delta_n$. If $u_n$ is
the resulting unitary, replace $R$ by $Ru_n$.

There is no circular choice: $G_n$ is chosen before $u_n$. At stage $n\geq2$, the product $R$ to be updated was the fixed opposite product when $L_{n-1}$
was formed. Hence 
\begin{equation*}
F\cup\{a_j:j\leq n\}\cup \mathrm{Ad}(R^*)\bigl(F\cup\{a_j:j\leq n\}\bigr)\subset L_{n-1}.
\end{equation*}
For $n=1$, the same inclusion follows from $R=1$ and $L_0=F\cup\{a_1\}$.
Therefore, for every $a\in F\cup\{a_j:j\leq n\}$, 
\begin{equation*}
\begin{split}
\left\|\mathrm{Ad}(Ru_n)(a)-\mathrm{Ad} R(a)\right\| &=\left\|\mathrm{Ad}
u_n(a)-a\right\|<\eta_n, \\
\left\|\mathrm{Ad}((Ru_n)^*)(a)-\mathrm{Ad}(R^*)(a)\right\| &=\left\|\mathrm{Ad} u_n^*(\mathrm{Ad}(R^*)(a))-\mathrm{Ad}(R^*)(a)\right\|<\eta_n.
\end{split}\end{equation*}
Here $\left\|\mathrm{Ad} u_n^*(b)-b\right\|=\left\|\mathrm{Ad}
u_n(b)-b\right\|$.

Write $P_n,Q_n$ for the products after stage $n$. Summability and density
now imply point-norm convergence of their adjoint actions and of the adjoint
actions of their inverses. Write 
\begin{equation*}
\mathrm{Ad} P_n\to\alpha,\quad \mathrm{Ad}(P_n^*)\to\widetilde\alpha, \qquad 
\mathrm{Ad} Q_n\to\beta,\quad \mathrm{Ad}(Q_n^*)\to\widetilde\beta.
\end{equation*}
For example, for every $a\in A$, 
\begin{equation*}
\begin{split}
\left\|\alpha(\widetilde\alpha(a))-a\right\| &\leq
\left\|\alpha(\widetilde\alpha(a))- \mathrm{Ad} P_n(\widetilde\alpha(a))\right\| \\
&\quad+\left\|\widetilde\alpha(a)-\mathrm{Ad}(P_n^*)(a)\right\|\longrightarrow0.
\end{split}\end{equation*}
The other compositions are identical, so $\widetilde\alpha=\alpha^{-1}$ and $\widetilde\beta=\beta^{-1}$; in particular $\alpha,\beta$ are automorphisms.

For $n\geq j$, the stage-$n$ state test gives 
\begin{equation*}
\left|\rho_0\mathrm{Ad} Q_n(a_j)-\rho_1\mathrm{Ad} P_n(a_j)\right|
<\delta_n<2^{-n}.
\end{equation*}
Passing to the limit and then using density gives $\rho_0\beta=\rho_1\alpha$. Thus $v=\alpha\beta^{-1}$ has $\rho_0v^{-1}=\rho_1$. Every control set
contains $F$. Applying the isometry $\alpha^{-1}$ gives $\left\|v(a)-a\right \|=\left\|\beta^{-1}(a)-\alpha^{-1}(a)\right\|$, and $\left\|\alpha^{-1}(a)-a\right\|=\left\|\alpha(a)-a\right\|$, so 
\begin{equation*}
\left\|v(a)-a\right\|\leq\left\|\beta^{-1}(a)-a\right\|
+\left\|\alpha(a)-a\right\| \leq2\sum_n\eta_n<\varepsilon\quad(a\in F).
\end{equation*}
Scaling handles finite sets outside the unit ball.
\end{proof}


\subsubsection{Completion and minimality}

\begin{proposition}\label{prop:cu-completion}
\uses{def:coprecompact}
 Let $A$ be simple, separable, unital, infinite
dimensional and nuclear, let $\varphi$ be a pure state, and put $G=\mathrm{Aut}(A)$ and $H=H_\varphi=\mathrm{Aut}(A,\varphi)$. The orbit map 
\begin{equation*}
G/H\longrightarrow P(A),\qquad gH\longmapsto\varphi\circ g^{-1}
\end{equation*}
is a uniform isomorphism from the quotient right uniformity onto the weak-*
subspace uniformity. Its compact completion is $S(A)$. In particular $H$ is
co-precompact.
\end{proposition}

\begin{proof}
\uses{lem:cu-glimm-avoid, thm:cu-alignment}
Pure-state homogeneity \cite{KishimotoOzawaSakai2003} makes the map bijective. The quotient
right uniformity is generated by pairs $(ugH,gH)$ for $u$ in a neighbourhood
of $1$. Their image states have the form $(\rho u^{-1},\rho)$, and 
\begin{equation*}
\left|\rho(u^{-1}(a))-\rho(a)\right|
\leq\left\|u^{-1}(a)-a\right\|=\left\|u(a)-a\right\|.
\end{equation*}
This proves uniform continuity. Theorem~\ref{thm:cu-alignment} proves uniform
continuity of the inverse, with tests independent of $\rho$. Pure states are
weak-* dense by Lemma~\ref{lem:cu-glimm-avoid}, applied in any essential
irreducible representation. Since the state space is compact, it is the
completion.
\end{proof}

\begin{corollary}[State-space criterion]\label{cor:cu-state-criterion}
\uses{prop:cu-completion}
 Under the hypotheses of Proposition~\ref{prop:cu-completion}, if $H_\varphi$ is extremely amenable and the action of $\mathrm{Aut}(A)$ on $S(A)$ is minimal, then 
\begin{equation*}
M(\mathrm{Aut}(A))\cong S(A).
\end{equation*}
\end{corollary}

\begin{proof}
\uses{lem:cu-completion-criterion, prop:cu-completion}
Apply Lemma~\ref{lem:cu-completion-criterion} to Proposition~\ref{prop:cu-completion}.
\end{proof}

\begin{proposition}\label{prop:cu-minimal}
 The action of $\mathrm{Aut}(\mathcal{O}_2)$ on $S(\mathcal{O}_2)$ is minimal and strongly proximal. Its pure-state orbit is
comeager.
\end{proposition}

\begin{proof}
Fix a source state $\rho$, a target pure state $\sigma$, and a finite set $F$
in the unit ball. Pure-state excision \cite[Proposition 2.2]{AkemannAndersonPedersen1986} first
gives a positive contraction $h$ of norm one, with tolerances chosen
sufficiently small, for which $\left\|hah-\sigma(a)h^2\right\|$ is small on $F$. For small $\gamma>0$, real rank zero \cite[Theorem 2.6(iii)]{BrownPedersen1991} gives a nonzero projection $p$ in the hereditary algebra
generated by $(h-(1-\gamma))_+$. Since $\left\|hp-p\right\|\leq\gamma$,
compression and the triangle inequality give 
\begin{equation*}
\left\|pap-\sigma(a)p\right\|<\epsilon\qquad(a\in F).
\end{equation*}
One may ensure properness by passing to a nonzero proper subprojection. For $L$ large, the $2^L$ Cuntz word-range projections sum to $1$; one of them, $z$, satisfies $\rho(z)\leq2^{-L}$. Put $q=1-z$. Nonzero projections in $\mathcal{O}_2$ are Murray--von Neumann equivalent \cite{Cuntz1981}. Applying
this to $p,q$ and to their complements gives a unitary $v$ with $vpv^*=q$.
The state $\rho^{\prime }=\rho\mathrm{Ad} v$ has $\rho^{\prime -L}$.
Cauchy--Schwarz, applied to the two off-diagonal corners of $a$, gives 
\begin{equation*}
\left|\rho^{\prime }(a)-\sigma(a)\right| \leq\epsilon+2\sqrt{2^{-L}}+2\,2^{-L}\qquad(a\in F).
\end{equation*}
Thus every orbit approaches every pure state. Pure-state density proves
minimality.

For strong proximality, let $\mu$ be a probability measure on $S(\mathcal{O}_2)$ and let $\bar\mu$ be its barycenter. By minimality, a net of translates
of $\bar\mu$ converges to a pure state $\sigma$. A cluster point of the
corresponding translated measures has barycenter $\sigma$; extremality
forces that measure to be $\delta_\sigma$. Finally the extreme points of a
metrizable compact convex set form a $G_\delta$. Here they are dense and
form one orbit.
\end{proof}

\begin{proof}[Completion of Theorem~\protect\ref{thm:cu-main}]
\proves{thm:cu-main}
\uses{cor:cu-h-ea, cor:cu-scalar-umf, lem:cu-completion-criterion, prop:cu-completion, prop:cu-minimal, thm:cu-w-ea}
Theorem~\ref{thm:cu-w-ea}, Corollary~\ref{cor:cu-h-ea}, and Corollary~\ref{cor:cu-scalar-umf} prove the first, second, and fourth assertions. For the
third, apply Lemma~\ref{lem:cu-completion-criterion}. Its hypotheses are
precisely Corollary~\ref{cor:cu-h-ea} and Propositions~\ref{prop:cu-completion}
and~\ref{prop:cu-minimal}.
\end{proof}

\begin{corollary}\label{cor:cu-furstenberg}
\uses{def:proximal}
 The state-space action is the universal minimal
strongly proximal flow, or Furstenberg boundary, of $\mathrm{Aut}(\mathcal{O}_2)$. Every minimal compact $\mathrm{Aut}(\mathcal{O}_2)$-flow is metrizable
and strongly proximal; in particular, $S(\mathcal{O}_2)$ is also the
universal minimal proximal flow.
\end{corollary}

\begin{proof}
\uses{prop:cu-minimal, thm:cu-main}
Theorem~\ref{thm:cu-main} identifies the universal minimal flow with $S(\mathcal{O}_2)$, and Proposition~\ref{prop:cu-minimal} makes this flow strongly
proximal. Every minimal flow is a factor of the universal minimal flow;
metrizability and strong proximality pass to factors.
\end{proof}

\begin{corollary}\label{cor:cu-radicals}
 The group $\mathrm{Aut}(\mathcal{O}_2)$ has trivial
Bohr compactification and no nontrivial closed normal amenable subgroup.
\end{corollary}

\begin{proof}
\uses{cor:cu-furstenberg}
A nontrivial compact group compactification of $\mathrm{Aut}(\mathcal{O}_2)$
would give a nontrivial minimal translation flow carrying Haar probability,
contradicting Corollary~\ref{cor:cu-furstenberg}. For the second assertion, let 
$N$ be closed, normal, and amenable. It fixes a probability $\mu$ on $S(\mathcal{O}_2)$. Normality makes every $G$-translate of $\mu$ again $N$-invariant. Strong proximality therefore gives an $N$-fixed state, and
normality plus minimality makes $N$ fix the whole state space. States
separate automorphisms, so $N$ is trivial.
\end{proof}

\begin{remark}\label{rem:cu-aut-nonamenable}
 The state-space flow has no invariant
probability. Such a probability would have an inner-invariant barycenter,
hence a tracial state on $\mathcal{O}_2$. The Cuntz relations would give $1=\tau(s_1s_1^*)+\tau(s_2s_2^*)=2$. Thus $\mathrm{Aut}(\mathcal{O}_2)$ is
nonamenable, although its UMF is metrizable.
\end{remark}


\subsection{Related Cuntz groups}

\label{sec:cu-related}

Throughout this section $A=\mathcal{O}_2$ and $D=A\otimes\mathcal{K}$. The
formulas specify the acting quotient as well as the compact space.


\subsubsection{Stable automorphisms and rigidity}

\begin{corollary}\label{cor:cu-stable-o2}
 Both $\mathrm{Aut}(D)$ and $\mathrm{Aut}(D^+)$, with
point-norm topology, are extremely amenable.
\end{corollary}

\begin{proof}
\uses{thm:cu-w-ea}
We show that $\mathrm{Ad} U_1(D)$ is dense in $\mathrm{Aut}(D)$. Let $\alpha\in\mathrm{Aut}(D)$. Approximate a given finite test set by a
canonical corner $pDp$. The projections $\alpha(p)$ and $p$ can be
conjugated by a unitary in $1+D$. Indeed approximate $\alpha(p)$ by a
projection in a finite canonical corner and first use the close-projection
unitary. Enlarge that corner to contain both projections with nonzero
complements. In this matrix algebra over $\mathcal{O}_2$, the projections
and their respective complements are equivalent; the implementing partial
isometries sum to a corner unitary. Extend it by the identity.

After this conjugation $\alpha$ preserves $pDp\cong\mathcal{O}_2$, where its
restriction is approximately inner by \cite{KirchbergPhillips2000}. Extending the resulting
corner unitary by $1-p$ approximates $\alpha$ on the original tests. Density
and Theorem~\ref{thm:cu-w-ea} imply extreme amenability. Every automorphism of $D^+$ preserves its unique nonzero proper ideal $D$, and every automorphism
of $D$ extends uniquely. Restriction is a point-norm topological isomorphism.
\end{proof}

\begin{proposition}\label{prop:cu-rigidity}
 The only $H_\varphi$-invariant state on $A$ is $\varphi $. In particular $N_{\mathrm{Aut}(A)}(H_\varphi)=H_\varphi$, and the
only continuous $\mathrm{Aut}(A)$-equivariant self-map of $S(A)$ is the
identity.
\end{proposition}

\begin{proof}
An invariant state $\rho$, restricted to $D_\varphi$, is invariant under all
its unitisation inner automorphisms, so its restriction is a bounded trace.
A stable algebra has no nonzero bounded trace: arbitrarily many orthogonal
equivalent copies of a positive element force its trace to vanish. Thus $\rho|_{D_\varphi}=0$. The support of $D_\varphi$ is $1-p_\varphi$, so the
normal extension of $\rho$ is supported on the minimal projection $p_\varphi$, forcing $\rho=\varphi$. If $g$ normalises $H_\varphi$, then $g\cdot\varphi$
is $H_\varphi$-fixed and hence equals $\varphi$. An equivariant self-map
likewise fixes $\varphi$, then the dense pure-state orbit, and therefore the
whole space.
\end{proof}

\begin{remark}\label{rem:cu-remark-3}
The norm groups $U(\mathcal{O}_2)$ and $\mathrm{PU}(\mathcal{O}_2)$ are
nonamenable, by the tracial-barycenter obstruction in Section~\ref{sec:cu-state}. Their UMFs are nonmetrizable: a Polish SIN group with metrizable UMF is
amenable \cite[Corollary 3.4]{BYMT2017}. These groups differ from both the
point-norm automorphism group and the stabilised unitisation group computed
above.
\end{remark}


\subsection{Stable UCT Kirchberg automorphism groups}

\label{sec:cu-kirchberg}

A Kirchberg algebra is nonzero, separable, nuclear, simple and purely
infinite. Here $B$ is stable and satisfies the universal coefficient theorem
(UCT). Write $\underline K(B)$ for total $K$-theory: the ordinary graded
groups, the groups with finite coefficients, and the coefficient and
Bockstein operations. Give 
\begin{equation*}
Q_B=\mathrm{Aut}_\Lambda(\underline K(B))
\end{equation*}
pointwise convergence on this countable discrete invariant. Preservation of
operations is closed in the corresponding permutation group, so $Q_B$ is
Polish. Coefficient information is necessary even for some finite $K$-groups.

\begin{theorem}\label{thm:cu-stable-general}
 The total-$K$-theory homomorphism is a continuous
open surjection 
\begin{equation*}
q_B:\mathrm{Aut}(B)\longrightarrow Q_B
\end{equation*}
whose kernel is extremely amenable. Consequently 
\begin{equation*}
M(\mathrm{Aut}(B))\cong q_B^*M(Q_B),
\end{equation*}
where $q_B^*$ denotes pullback of the action.
\end{theorem}

\begin{proof}
\uses{cor:cu-cuntz-stable, lem:cu-kernel, lem:cu-strict-ea}
Let 
\begin{equation*}
N_B=\overline{\{\mathrm{Ad} u:u\in U(M(B))\}}^{\,\text{point-norm}}.
\end{equation*}
Choose a nonzero projection $p\in B$. It is full, and Brown's theorem \cite[Theorem 2.8]{Brown1977} gives $B\cong B_0\otimes\mathcal{K}$ for the unital
Kirchberg algebra $B_0=pBp$. The adjoint map $\mathrm{Ad}\colon U(M(B))\to \mathrm{Aut}(B)$ is continuous for the strict topology on the source and the
point-norm topology on the target: strict convergence controls $(u_i-u)a$
and $(u_i-u)a^*$, which bound the two terms in $\mathrm{Ad} u_i(a)-\mathrm{Ad} u(a)$. Lemma~\ref{lem:cu-strict-ea} below makes $U(M(B))$ extremely amenable,
and dense-image permanence makes $N_B$ extremely amenable.

With the standard convention for stable algebras, approximate unitary
equivalence is implemented by multiplier unitaries. Thus the group denoted $\overline{\mathrm{Inn}}(B)$ in the classification literature is precisely $N_B$. Kirchberg--Phillips classification together with the universal
multicoefficient theorem \cite{KirchbergPhillips2000,Phillips2000,RosenbergSchochet1987,DadarlatLoring1996,Dadarlat2005} give the exact
sequence 
\begin{equation}  \label{eq:cu-classification-sequence}
1\longrightarrow N_B\longrightarrow\mathrm{Aut}(B) \xrightarrow{q_B}Q_B\longrightarrow1.
\end{equation}
Surjectivity can be read directly from those ingredients. The universal
multicoefficient theorem lifts any $\theta\in Q_B$ to $x\in KK(B,B)$, and
invertibility of $\theta$ on total $K$-theory makes $x$ invertible. Via $B\cong B_0\otimes\mathcal{K}$ and Morita invariance, regard $x$ as a class
in $KK(B_0,B_0)$. Stable Kirchberg classification \cite[Theorem 4.2.1]{Phillips2000} then realises the class by an isomorphism of stabilisations,
hence by an automorphism of $B_0\otimes\mathcal{K}\cong B$. The kernel
assertion in \eqref{eq:cu-classification-sequence} says exactly that trivial
total $K$-theory class is equivalent to approximate unitary equivalence to
the identity. Continuity can be checked on representatives: each coefficient 
$K$-class is represented by a projection or unitary in a fixed coefficient
algebra and amplification, and a sufficiently small perturbation leaves its
class unchanged. Dadarlat identifies the corresponding quotient topology on $KK/\overline 0$ with pointwise convergence on total $K$-theory \cite[Theorem
4.1(b)]{Dadarlat2005}. A continuous surjective homomorphism between Polish
groups is open. Apply Lemma~\ref{lem:cu-kernel} to this topological exact
sequence. The group $Q_B$ is discrete precisely when $\underline K(B)$ is
finitely generated. If in addition $Q_B$ is finite, then $M(\mathrm{Aut}(B))\cong Q_B$ as a translation flow; this is the case for the stabilised
Cuntz algebras in Corollary~\ref{cor:cu-cuntz-stable}. If $Q_B$ is infinite
discrete, $M(\mathrm{Aut}(B))$ is the pullback of the (typically
nonmetrizable) universal minimal flow of that discrete group, not $Q_B$
itself. For infinitely generated invariants, $Q_B$ need not be discrete and
the formula remains a pullback of $M(Q_B)$.
\end{proof}

\begin{lemma}[Strict multiplier unitaries]\label{lem:cu-strict-ea}
\uses{def:ea}
 Let $A$ be a unital separable simple nuclear C*-algebra and put $B=A\otimes\mathcal{K}$. The group $U(M(B))$, equipped with
the strict topology, is extremely amenable. In particular this applies when $B$ is a stable Kirchberg algebra.
\end{lemma}

\begin{proof}
The group is Polish. Nuclearity of $A$ is equivalent to the AF-property of $M(B)$ in the strict topology \cite[Theorem~3.3]{NgRuiz2012}: there is an
increasing sequence of finite-dimensional unital C*-subalgebras $F_n\subset M(B)$ whose union is strictly dense. A Kaplansky-type density
statement in the same paper \cite[Proposition~4.1(vi)]{NgRuiz2012} then makes $\bigcup_n U(F_n)$ strictly dense in $U(M(B))$.

Stability supplies a coherent amplification that keeps the sequence
increasing. Identify $B\cong M_n(B)$ and therefore $M(B)\cong M_n(M(B))$.
Transport $F_n$ to $F_n^{\prime }:=F_n\otimes 1_{M_n}$ under this
isomorphism. Every matrix summand of $F_n^{\prime }$ then has size at least $n$, and one may choose the identifications so that $F_n^{\prime }\subset
F_{n+1}^{\prime }$. Write $U(F_n^{\prime })$ again as $U(F_n)$. Write $F_n\cong\bigoplus_iM_{k_i}$, so that $U(F_n)=\prod_iU(k_i)$, and let $\left\|\cdot\right\|_{\mathrm{HS}}$ denote the unnormalized Hilbert--Schmidt
norm of this finite-dimensional algebra. For $b\in B$ the strict seminorm
obeys 
\begin{equation*}
\left\|(u-v)b\right\|\leq\left\|b\right\|\,\left\|u-v\right\|
\leq\left\|b\right\|\,\left\|u-v\right\|_{\mathrm{HS}},
\end{equation*}
so every strict entourage of $U(F_n)$ contains a Hilbert--Schmidt ball of
fixed radius, and concentration in a coarser metric is implied by
concentration in a finer one. On $U(k)$, Gromov--Milman concentration in the
unnormalized Hilbert--Schmidt metric \cite{GromovMilman1983} bounds the deviation of a $1$-Lipschitz function from its mean by $2e^{-ckt^2}$, and tensorization gives
the same bound on the product with $k=\min_ik_i$. The amplification above
forces $\min_ik_i\geq n$. Hence $\bigl(U(F_n)\bigr)_n$ is a L\'evy family
for the ambient strict entourages. The union $\bigcup_n U(F_n)$ is still
strictly dense, so $U(M(B))$ is a L\'evy group in the sense of
Giordano--Pestov \cite[Section~2]{GiordanoPestov2007} and is therefore extremely amenable.
\end{proof}

\begin{remark}\label{rem:cu-remark-4}
\uses{cor:cu-stable-o2, lem:cu-strict-ea, thm:cu-main}
The strict topology is essential. This argument concerns point-norm
automorphisms and does not prove amenability of norm $U(B^+)$. Ku\v{c}erovsk\'y--Ng \cite[Theorem~3.11]{KucerovskyNg2009} obtain fixed points on compact
metric spaces for the strict unitary group of a stable AH-algebra; that
statement is close to extreme amenability but is restricted to the AH case.
The original citation \cite{KucerovskyNg2010} established nuclearity criteria via
absorbing extensions and does not by itself supply Lemma~\ref{lem:cu-strict-ea}. Ng--Ruiz prove the corresponding structural statement for the projective
group, namely that $U(M(A\otimes\mathcal{K}))/\mathbb{T}$ is the closure of
an increasing union of compact subgroups \cite[Theorem 4.2(ii)]{NgRuiz2012}; the
concentration estimate is what is added here. The lemma is not used for
Theorem~\ref{thm:cu-main}: extreme amenability of $\mathrm{Aut}(D)$ is already
Corollary~\ref{cor:cu-stable-o2}. It is used only for the Kirchberg table in
this section.
\end{remark}


\subsubsection{Ordinary and total \(K\)-theory}

Let $G_i=K_i(B)$, and set 
\begin{equation*}
E_B=\mathrm{Ext}^1_\mathbb{Z}(G_0,G_1)\oplus\mathrm{Ext}^1_\mathbb{Z}(G_1,G_0),\qquad P_B=\mathrm{Pext}^1_\mathbb{Z}(G_0,G_1)\oplus\mathrm{Pext}^1_\mathbb{Z}(G_1,G_0).
\end{equation*}
Put $R_B=\mathrm{End}_\Lambda(\underline K(B))$. The UCT maps the Kasparov
endomorphism ring onto $\mathrm{End}(G_0)\times\mathrm{End}(G_1)$ with
kernel $E_B$, whereas the multicoefficient theorem maps the same ring onto $R_B$ with kernel $P_B\subset E_B$. Passing between these two quotients gives
the exact sequence of rings 
\begin{equation*}
0\longrightarrow E_B/P_B\longrightarrow R_B \longrightarrow\mathrm{End}(G_0)\times\mathrm{End}(G_1)\longrightarrow0.
\end{equation*}
The kernel is square-zero: products of filtration-one UCT classes land in $\mathrm{Ext}^2_\mathbb{Z}=0$. Units lift modulo a square-zero ideal, and the
units in the kernel are exactly $1+E_B/P_B$, with group law given by
addition. Passing to units, with $e\mapsto1+e$, therefore gives the group
exact sequence 
\begin{equation}  \label{eq:cu-total-extension}
1\longrightarrow(E_B/P_B,+)\longrightarrow Q_B \longrightarrow\mathrm{Aut}(G_0)\times\mathrm{Aut}(G_1)\longrightarrow1.
\end{equation}
Quotienting by pure extensions gives the total-$K$-theory ring \cite{RosenbergSchochet1987,DadarlatLoring1996,Dadarlat2005}. No splitting of the displayed sequence is asserted.
Equation \eqref{eq:cu-total-extension} is used algebraically; the topology on $Q_B$ remains pointwise convergence.

\begin{corollary}\label{cor:cu-k1zero}
\uses{thm:ellis}
 If $K_1(B)=0$, then, topologically, 
\begin{equation*}
Q_B\cong\mathrm{Aut}(K_0(B)),\qquad M(\mathrm{Aut}(B))\cong M(\mathrm{Aut}(K_0(B))).
\end{equation*}
The right-hand group carries pointwise convergence on $K_0(B)$. The degrees
may be interchanged.
\end{corollary}

\begin{proof}
Both cross-degree extension terms vanish. Restriction is a continuous
bijection of Polish groups and hence a topological isomorphism.
\end{proof}

\begin{corollary}\label{cor:cu-cuntz-stable}
\uses{thm:ellis}
 For $2\leq n<\infty$, 
\begin{equation*}
M(\mathrm{Aut}(\mathcal{O}_n\otimes\mathcal{K})) \cong(\mathbb{Z}/(n-1)\mathbb{Z})^\times.
\end{equation*}
For $\mathcal{O}_\infty$, the stable automorphism UMF is the two-point group 
$\{1,-1\}$. These are translation actions through $K_0$. For $n=2$ the
quotient is trivial.
\end{corollary}

\begin{proof}
Use $K_0(\mathcal{O}_n)=\mathbb{Z}/(n-1)\mathbb{Z}$, $K_0(\mathcal{O}_\infty)=\mathbb{Z}$, and $K_1=0$ \cite{Cuntz1981}. Their automorphism groups
are finite and equal to their own universal minimal flows. Table~\ref{eq:cu-cuntz} lists small cases.
\end{proof}

\begin{table}[ht]
\centering
\begin{tabular}{lll}
\hline Stable algebra $B$ & $K_0(B)$ & $M(\mathrm{Aut}(B))$ \\ 
\hline $\mathcal{O}_2\otimes\mathcal{K}$ & $0$ & singleton \\ 
$\mathcal{O}_3\otimes\mathcal{K}$ & $\mathbb{Z}/2$ & singleton \\ 
$\mathcal{O}_4\otimes\mathcal{K}$ & $\mathbb{Z}/3$ & $C_2$ \\ 
$\mathcal{O}_6\otimes\mathcal{K}$ & $\mathbb{Z}/5$ & $C_4$ \\ 
$\mathcal{O}_9\otimes\mathcal{K}$ & $\mathbb{Z}/8$ & $C_2\times C_2$ \\ 
$\mathcal{O}_\infty\otimes\mathcal{K}$ & $\mathbb{Z}$ & $C_2$ \\ 
\hline &  & 
\end{tabular}\caption{Stable Cuntz algebras and the universal minimal flows of their
automorphism groups, from Corollary~\protect\ref{cor:cu-k1zero}.}
\label{eq:cu-cuntz}
\end{table}

\begin{example}\label{ex:cu-example-1}
The range theorem \cite[Theorem 4.2.5]{Phillips2000}, followed by stabilisation,
realises arbitrary countable graded abelian $K$-groups by stable UCT
Kirchberg algebras. Let $p$ be prime and $d\geq1$. If $G_0=(\mathbb{Z}/p)^d$
and $G_1=0$, the UMF is the finite translation group $\mathrm{GL}_d(\mathbb{F}_p)$. If $G_0=\mathbb{Z}[1/p]$ and $G_1=0$, then 
\begin{equation*}
Q_B=\{\pm p^k:k\in\mathbb{Z}\}\cong C_2\times\mathbb{Z}.
\end{equation*}
It is discrete because fixing $1$ forces the identity. Its UMF is $C_2\times
M(\mathbb{Z})$, which is nonmetrizable.
\end{example}

\begin{example}\label{ex:cu-coefficient}
 If $G_0=G_1=\mathbb{Z}/2$, the ordinary graded
automorphism group is trivial. However $P_B=0$, since pure extensions of
finitely generated groups split, and $E_B=(\mathbb{Z}/2)^2$. Thus $Q_B=(\mathbb{Z}/2)^2$ and the automorphism UMF is its four-point translation
action. Replacing total $K$-theory by the ordinary graded invariant would
give a wrong answer.
\end{example}

\begin{proposition}\label{prop:cu-finite-k}
\uses{thm:ellis}
 Suppose $G_0,G_1$ are finitely generated. Then $Q_B$
is countable discrete, and finite exactly when both free ranks are at most
one. In the finite case, 
\begin{equation*}
|M(\mathrm{Aut}(B))| =|\mathrm{Aut}(G_0)|\,|\mathrm{Aut}(G_1)|\, |\mathrm{Ext}^1_\mathbb{Z}(G_0,G_1)|\,|\mathrm{Ext}^1_\mathbb{Z}(G_1,G_0)|.
\end{equation*}
If either free rank is at least two, $M(\mathrm{Aut}(B))$ is nonmetrizable.
\end{proposition}

\begin{proof}
Pure extensions vanish, and both extension groups are finite. Equation \eqref{eq:cu-total-extension} makes $Q_B$ countable. A countable Polish group
is discrete by Baire. For $\mathbb{Z}^r\oplus T$, with $T$ finite, the
automorphism group maps onto $\mathrm{GL}_r(\mathbb{Z})$ with finite kernel.
It is finite exactly when $r\leq1$, proving the criterion and the
cardinality formula. Otherwise $Q_B$ is infinite discrete, hence locally
compact and noncompact, and therefore has nonmetrizable universal minimal
flow by Veech's theorem together with \cite[Corollary 3.9 and the remark
following it]{BYMT2017}; for locally compact groups this is the theorem of
Kechris--Pestov--Todor\v{c}evi\'c \cite{KPT2005}.
\end{proof}

For general countable $K$-groups, Theorem~\ref{thm:cu-stable-general} is a
reduction to $Q_B$; it does not assert that every such group's UMF has a
known geometric model.

\begin{example}[A natural-order flow]\label{ex:cu-vector-space}
 For a prime $p$, let $V_p=\bigoplus_{j=1}^{\infty}\mathbb{Z}/p\mathbb{Z}$, viewed as the countable-dimensional vector space
over $\mathbb{F}_p$, and use the range theorem to choose a stable UCT
Kirchberg algebra $B_p$ with $K_0(B_p)=V_p$ and $K_1(B_p)=0$. Then 
\begin{equation*}
M(\mathrm{Aut}(B_p))\cong M(\mathrm{GL}(V_p)) \cong\mathrm{NO}(V_p),
\end{equation*}
where $\mathrm{NO}(V_p)$ is the space of those linear orders on $V_p$ whose
restriction to every finite-dimensional subspace is natural, that is,
antilexicographic with respect to some ordered basis. It is a closed, hence
compact, subspace of $2^{V_p\times V_p}$. The last identification is the
vector-space computation in the KPT correspondence \cite[Theorem 8.2(ii)]{KPT2005}. This gives an infinite, metrizable, nontrivial UMF for a stable
Kirchberg automorphism group.
\end{example}

\begin{example}[Compact $p$-adic flows]\label{ex:cu-prufer}
 For a prime $p$ and $d\geq1$, let $C_{p^\infty}$ be the
Pr\"ufer $p$-group and let $B_{p,d}$ be a stable UCT Kirchberg algebra with $K_0(B_{p,d})=C_{p^\infty}^{\,d}$ and $K_1(B_{p,d})=0$. Then 
\begin{equation*}
\mathrm{Aut}(K_0(B_{p,d}))\cong\mathrm{GL}_d(\mathbb{Z}_p)
\end{equation*}
with its compact $p$-adic topology: fixing the subgroup of elements of order
at most $p^n$ is congruence modulo $p^n$. Indeed, restriction to the $p^n$-torsion subgroups identifies $\mathrm{End}(C_{p^\infty})$ with $\varprojlim_n\mathbb{Z}/p^n\mathbb{Z}=\mathbb{Z}_p$; taking matrices and
units gives the displayed topological group. Hence 
\begin{equation*}
M(\mathrm{Aut}(B_{p,d}))\cong\mathrm{GL}_d(\mathbb{Z}_p)
\end{equation*}
as a translation flow. For $d=1$ this is $\mathbb{Z}_p^\times$.
\end{example}

\subsection{Proofs of statements in Section~\ref{sec:operator-algebras}}\label{sec:cu-encyclopaedia-proofs}

\begin{proof}[Proof of Theorem~\ref{thm:lupini-cuntz}]
\proves{thm:lupini-cuntz}
\uses{thm:cu-main}
This is Theorem~\ref{thm:cu-main}.
\end{proof}

\begin{proof}[Proof of Example~\ref{ex:lupini-stable-cuntz}]
\proves{ex:lupini-stable-cuntz}
\uses{cor:cu-k1zero, cor:cu-cuntz-stable}
By Corollary~\ref{cor:cu-k1zero} the quotient is $\mathrm{Aut}(K_0(A))$ when
$K_1(A)=0$, and Corollary~\ref{cor:cu-cuntz-stable} computes the universal
minimal flows: $(\mathbb{Z}/(n-1))^\times$ for $2\leq n<\infty$, and
$\{1,-1\}$ for $\mathcal{O}_\infty\otimes\mathcal{K}$. For $n=2$ and $n=3$ the
group $(\mathbb{Z}/(n-1))^\times$ is trivial.
\end{proof}
